%! Author = taj
%! Date = 7/22/25

% Preamble
\documentclass[11pt]{article}

% Packages
\usepackage{amsmath}
\usepackage{graphicx}
\usepackage{float}
\usepackage{upgreek}

\title{9. Precesija}
\author{Taj Šobot}

% Document
\begin{document}
    \maketitle
    \flushleft
    Pri tem eksperimentu smo merili precesijo kolesa in njen vstrajnostni moment.

    \section{Meritve}
    Masa kolesa in uteži:
    \begin{align*}
        m_k = (937 \pm 1) g
        m_u = (51 \pm 1) g
    \end{align*}

    Skica eksperimenta:
    \begin{figure}[H]
        \centering
        \includegraphics[width=0.9\textwidth]{img}
        \caption{Skica eksperimenta}
        \label{fig:1}
    \end{figure}
    Meritve precesije smo naredili v položaju prikazanem na skici:
    Merili smo začetno in končno frekvenco, ter čas ene precesije.
    Tabela:


    \begin{table}[h!]
        \centering
        \begin{tabular}{|c|c|c|c|}
            \hline
            n & $\upnu_z [rpm]$ & $\upnu_k [rpm]$ & $t_0 [s]$ \\
            \hline
            1  & 1023 & 761 & 7.16 \\
            2  & 977  & 951 & 7.07 \\
            3  & 959  & 813 & 7.66 \\
            4  & 1061 & 863 & 7.67 \\
            5  & 989  & 796 & 6.22 \\
            6  & 585  & 748 & 6.34 \\
            7  & 970  & 799 & 7.04 \\
            8  & 799  & 685 & 6.08 \\
            9  & 1068 & 807 & 7.17 \\
            10 & 1226 & 827 & 7.40 \\
            \hline
        \end{tabular}
        \caption{tabela}
        \label{tab:1}
    \end{table}
    Kolo smo potem dali še na stojalo in mu dodali utež.
    To smo potem zanihali (da lahko kasneje izračunamo vstrajnostni moment).
    Nihajni čas(5 nihajev):
    \[
        T_{0 za 5 nihajev} = (16.13 \pm 0.1) s
    \]
    Razdalja klade od sredine:
    \[
        r' = (19.5 \pm 0.5) cm
    \]

    \section{Izračuni in rezultati}
    Začnimo z določanjem vstrajnostnega momenta.
    Nihajni čas enega nihaja je bil:
    \[
        t_0 = (3.23 \pm 0.02)
    \]
    Vstrajnostni moment klade in kolesa je:
    \[
        J_s = J_k + m_u r'^2
    \]
    Vemo tudi da je:
    \[
        t_0 = 2\pi \sqrt[2]{\frac{J_s}{m_u g r'}}
    \]
    Obrnemo...
    \[
        J_s = \frac{t_0^2}{4 \pi^2} m_u g r'
    \]
    Združimo in dobimo:
    \[
        J_{k} = \frac{t_0^2}{4 \pi^2} m_u g r' - m_u r'^2
    \]
    Vstavimo in dobimo:
    \[
        J_k = 0.024(1 \pm 0.07) kg m^2
    \]
    To vrednost nato uporabimo v enačbi:
    \[
        \omega_p = \frac{m g r}{J-k \omega}
    \]
    kjer je:
    \[
        \overline{\omega}= 2 \pi (\frac{\upnu_z + \upnu_k}{2})
    \]
    Vstavimo vrednosti:
    \begin{align*}
        \overline{\upnu_z} = (923 \pm 32)rpm = (3.9 \pm 1) 1/s \\
        \overline{\upnu_k} = (738 \pm 21)rpm = (3.1 \pm 1) 1/s \\
        \overline{t_0} = (6.4 \pm 0.2)s
    \end{align*}
    In dobimo:
    \[
        \omega_p = (0.99  \pm 0.07) 1/s
    \]
    Kotno hitrost precesije lahko izračunamo tudi na enostavni način:
    \[
        \omega_p = \frac{2 \pi}{t_{0_p}} =(0.98 \pm 0.03) 1/s
    \]
    Ta rezultat zgleda pravilno.

    \section{Odgovori na vprašanja}
    \textit{Izpeljite enačbo za kotno hitrost precesije ob predpostavki, da je navor trenja v ležaju kolesa
    linearno odvisen od frekvence vrtenja kolesa.}
    \\
    Izpeljava:
    \[
        M_t = -k \upnu_t = -\frac{k}{2 \pi} \omega(t)
    \]
    M izrazimo drugače:
    \[
        J x = -\frac{k}{2 \pi} \omega(t)
    \]
    X izrazimo z...
    \[
        J \frac{d\omega}{dt} = -\frac{k}{2 \pi} \omega(t)
    \]
    Integriramo...
    \[
        \omega(t) = \omega_0 e^{\frac{k t}{2 \pi J}}
    \]
    Frekvenca...
    \[
        \upnu(t) = \upnu_0  e^{\frac{k t}{2 \pi J}}
    \]
    To vstavimo v enačbo za povprečno kotno hitrost (solata) in dobimo:
    \[
        \omega_p = \frac{4 m g r}{J \pi \upnu_z (1 + e^{\frac{k t}{2 \pi J}} )}
    \]


\end{document}